\section{The corrected annular limit on the full norm-one-vanishing test class}
\label{sec:local-extension}

The condition in task~1 is $h(1)=0$, not that the support of
$h$ be separated from $1$. The following extension closes that
local distinction. It does not treat nonzero $h(1)$, for which
a logarithmic subtraction is required, and it does not replace
the global arithmetic transport.

\paragraph{MA-LA-007 (all compact tests vanishing at norm one).}
In MA-LA-003 let instead $h\in C_c^\infty(\mathbb R_{>0})$
with $h(1)=0$. The exact even-sector trace (MA-LA.5) still holds.
Writing $L=XY$ and $r=Y\epsilon$, as there, one has
\[
 \operatorname{Tr}_{\rm even}T_{X,Y,\epsilon}
 \longrightarrow A_0(h)
 \quad\hbox{if }L\to\infty,\ r\to0,
 \qquad
 \operatorname{Tr}_{\rm even}T_{X,Y,\epsilon}
 \longrightarrow0
 \quad\hbox{if }r\to\infty,
 \tag{MA-LA.19}
\]
where
\[
 \begin{split}
 A_0(h)
 &=\frac12\int_0^\infty h(t)t^{-1/2}
       \left(\frac1{|t^{-1}-1|}+\frac1{t^{-1}+1}\right)\frac{dt}{t}\\
 &=\int_0^\infty\frac{e^{-x/2}}{1-e^{-2x}}
                 \bigl(h(e^x)+h(e^{-x})\bigr)\,dx.
 \end{split}
 \tag{MA-LA.20}
\]
The integrals are absolutely convergent. For fixed smooth
profiles $A,B$, the first error is $O_h(L^{-1}+r)$ for
$L\ge1$, $0<r\le1$, and the absolute trace in the second
regime is $O_h(r^{-1})$ for $r\ge1$.

For the thin boundary smoothings of MA-LA-004, with
$\delta<\epsilon/4$, the sufficient-rate error becomes
\[
 O_h\!\left(\frac{1+\log L}{L}+r+\delta(X+Y)\right).
 \tag{MA-LA.21}
\]
When $L,r\to\infty$ their absolute trace is bounded by
\[
 O_h\!\left(\frac{1+\log L}{L}
             +\frac{1+\log r}{r}+\delta(X+Y)\right).
 \tag{MA-LA.22}
\]
Thus both limiting paths in (MA-LA.12) remain valid on the full
requested test class, with the same two different limits whenever
$A_0(h)\ne0$.

\begin{proof}
The kernel in (MA-LA.7) is Schwartz for every fixed compact
support of $h$ inside $(0,\infty)$, whether or not that support
meets $1$. The trace-class and reflection arguments of
MA-LA-003 therefore give the same exact formula. Put
$c_-(t)=t^{-1}-1$ and $c_+(t)=t^{-1}+1$. The support is
contained in $[a,b]$ for some $0<a<b<\infty$.
For $t\ne1$, the absolute value of the inner integral
in (MA-LA.8) is at most
\[
 \|\check B\|_{L^1}
      \left(\frac1{|c_-(t)|}+\frac1{c_+(t)}\right),
\]
because the multiplier has absolute value at most one.
The factor $|h(t)|/|c_-(t)|$ is bounded near $1$:
$h(t)=O(|t-1|)$, while
$|c_-(t)|=|t-1|/t$. The other denominator stays bounded
away from zero throughout the compact support. Including
the bounded $t^{-1/2}$ and $dt/t$ factors, this is an
integrable dominating function in $t$. The point $t=1$
has measure zero and $h(1)=0$.

For each fixed $t\ne1$ the limiting inner integral has already
been computed in MA-LA-003. Dominated convergence now gives
both limits. For $t>1$, the first density in (MA-LA.20) is
$\sqrt t/(t^2-1)\,dt$; for $t<1$ it is
$t^{-1/2}/(1-t^2)\,dt$. Substituting $t=e^x$ and
$t=e^{-x}$, respectively, gives the second expression and
its absolute convergence.

For the quantitative bounds, the part removed by the outer
profile is at most
\[
 C_N|c|^{-1}(1+|c|L)^{-N}
\]
for either nonzero $c=c_\pm(t)$, by the Schwartz tail of
$\check B$. In a fixed neighborhood of $1$, multiply by
$|h(t)|$, cancel the factor $|c_-|^{-1}$, and use that
$dc_-/dt=-t^{-2}$ is bounded away from zero. For $N>1$,
\[
 \int_{|u|\le C}(1+|u|L)^{-N}\,du=O(L^{-1}).
\]
Away from that neighborhood the bound is $O_N(L^{-N})$;
the plus sign is harmless everywhere. The inner part removed
when $r\le1$ is supported on $|z|\le Cr$ and bounded by
$C\|\check B\|_\infty$, so its integral is $O(r)$,
uniformly in $t$. This proves the first error estimate.
When $r\ge1$, the remaining multiplier vanishes on
$|z|\le cr$. Replace $L$ by $r$ in the Schwartz-tail
argument just given to obtain $O(r^{-1})$.

For thin smoothings the error comparison with the sharp
annulus and sharp frequency interval is still
$O_h(\delta(X+Y))$: its proof used only bounded $t$ factors
and compact support, never separation of $c_-$ from zero.
For $u>0$ the sine integral satisfies
\[
 \left|\operatorname{Si}(u)-\frac\pi2\right|
       \le C\min(1,u^{-1}),\qquad
 |\operatorname{Si}(u)|\le C\min(u,1).
\]
The first inequality follows from integration by parts when
$u\ge1$, with boundedness for smaller $u$; the second is
immediate near zero and follows from boundedness otherwise.
In (MA-LA.11), the outer-end error is thus bounded by a
constant times
\[
 \int_{|v|\le C}\min(1,(|v|L)^{-1})\,dv
       =O((1+\log L)/L),
\]
using the same cancellation $|h|/|c_-|=O(1)$ and change
of variable. Away from $1$ the error is smaller. If $r\le1$,
the inner sine integral gives $|\operatorname{Si}(2\pi|c|r)|
\le C|c|r$, which yields $O_h(r)$. If $r\ge1$, apply
the outer-end estimate also with $L$ replaced by $r$.
Together with the smoothing comparison, these prove
(MA-LA.21)--(MA-LA.22).
\end{proof}

\paragraph{Scope, inputs, and programme placement.}
Every step uses the actual real dilation kernel and ordinary
Fourier analysis, with no zero-location premise. The sufficient
rate remains $\lambda'/\lambda''\to0$ in the notation of the
footnote on Masullo p.~514 and the proposed correction to
(12.4.1.23), p.~518. For the requested norm-one-vanishing
tests, this now proves the local annular repair even when their
support meets norm one. It closes the specific local question
Q-MA-LA-ENDPOINT; the nonzero-value logarithmic subtraction
and general-number-field cutoff estimates are different
questions and are not claimed.

In the adelic-trace interpretation this is the real-place
distribution on the full test class of GRH\_NIGHT MT-M30.
It refines the source-kernel side of ALPOGE\_BRIDGES
F-SM-20/21 without supplying the global prime transport.
The corresponding cote~1 issue (scans~079--080 and 111--112,
via BR-TR-01/02) is continuity of regularized trace under the
specified limiting process. No algebraic correspondence,
motivic weight statement, or new EGA/SGA identification is
asserted by this analytic estimate.
